\section{RELATED WORK}\label{related-work}

Single-machine loss-grid optimization intersects loss-landscape
visualization, adaptive loss-grid sampling, heterogeneous
CPU/GPU scheduling, and autotune-and-cache patterns.

\subsection{Loss-landscape visualization and
analytics}\label{rw-losslandscape}

Goodfellow et al.~\cite{goodfellow2015qualitative} introduced
1D linear interpolation between parameter points to
qualitatively characterize neural-network optimization
trajectories. Li et al.~\cite{li2018losslandscape} extended this
to a 2D loss grid with filter-normalized perturbation directions,
the algorithm taken here as the baseline; the official
implementation~\cite{goldstein_loss_landscape_repo} runs it under
MPI on a multi-node cluster. LossLens~\cite{xie2025losslens,losslens_repo}
integrated this construct into an interactive analytics
workflow over model modes/checkpoints.

\subsection{Adaptive and progressive loss-grid
sampling}\label{rw-adaptive-sampling}

Another line of work replaces the canonical full grid with
an adaptive sampler that evaluates loss only at points predicted
to carry visual information. Prakash et al.
apply this to loss-landscape visualization directly, replacing
regular $\alpha,\beta$ sampling with an adaptive algorithm
focused on ``interesting'' regions and comparing it against a
GradVis-style baseline and a vectorized uniform
sampler~\cite{prakash2024adaptive}. Their reported outcome is
mixed: adaptive sampling beats the GradVis-style baseline but
loses to the vectorized uniform sampler due to selection
overhead, and their user-preference comparison is
inconclusive~\cite{prakash2024adaptive}.

\subsection{Heterogeneous compute and ML inference
throughput}\label{rw-scheduling}

Task-parallel heterogeneous scheduling assigns whole
independent tasks across devices. Dynamic self-scheduling,
introduced by Polychronopoulos and
Kuck~\cite{polychronopoulos1987gss}, is the simplest variant:
idle workers poll a shared queue and claim one task at a time,
requiring no per-device performance model. Richer policies in
the same family include HEFT-style cost-model scheduling and
priority-based work-stealing, exemplified by
StarPU~\cite{augonnet2011starpu}; these address load
imbalance and per-task cost variance. Gallet and
Gowanlock~\cite{gallet2021heterogeneous} identify the
CPU/GPU throughput ratio as an applicability variable for
heterogeneous task scheduling, demonstrated on a CPU--GPU
epsilon grid join: static and dynamic partitioning yield a
speedup over GPU-only execution only above a workload-dependent
ratio threshold, and regress below it.

An orthogonal approach splits a single computation across
devices, with per-device stages exchanging intermediate data
along a pipeline. OpenVINO's HETERO execution mode is a
production example at the intersection of heterogeneous
compute and ML inference~\cite{openvino_hetero}: with
pipeline-parallel \texttt{model\_distribution\_policy}, the
inference model is divided into stages, each stage is assigned
to a device, and the output of one stage is fed as input to
the next.

Production ML inference systems also pursue throughput
through request-level batching. OpenVINO's automatic
batching~\cite{openvino_auto_batching,openvino_auto_batching_src}
groups concurrent batch-1 inference requests at runtime to
improve device utilization on supported devices and
throughput-oriented modes. At implementation level,
auto-batching operates on the input/request batch dimension
of one fixed compiled model: requests are packed into slices
of a larger input/output tensor and executed by a reshaped
batched model.

\subsection{Autotune-and-cache patterns}\label{rw-autotuning}

ATLAS~\cite{whaley1998atlas} introduces an autotune-and-cache
pattern: empirically tune the inner loops of BLAS routines
(GEMM, DAXPY, etc.) for the target machine's cache and
register hierarchy; cache the selected configuration by
(machine, BLAS-op) signature; reuse it across every
subsequent matching call. The pattern is four-part: tune
once per machine, cache by signature, reuse across matching
invocations, per-machine scoping. OpenTuner~\cite{ansel2014opentuner}
canonized bounded autotuning search with explicit stopping
criteria as the discipline-level departure from ATLAS-style
exhaustive search.

\subsection{Where this thesis sits}\label{rw-placement}

Prior work provides the core
algorithm~\cite{li2018losslandscape,goldstein_loss_landscape_repo}
and its interactive checkpoint-comparison use
case~\cite{xie2025losslens,losslens_repo}, the dynamic
self-scheduling
primitive~\cite{polychronopoulos1987gss} and the broader
heterocompute-scheduler family~\cite{augonnet2011starpu},
the throughput-ratio applicability
variable~\cite{gallet2021heterogeneous} from a different
workload, the production model-parallel and request-batching
counterparts in ML
inference~\cite{openvino_hetero,openvino_auto_batching}, the
autotune-and-cache pattern~\cite{whaley1998atlas} with its
bounded-search discipline~\cite{ansel2014opentuner}, and the
adjacent adaptive-sampling
direction~\cite{prakash2024adaptive}.

The gap is the combination of workload, deployment target, and
optimization question. Existing loss-landscape work defines the
canonical grid and motivates interactive checkpoint comparison,
but does not evaluate commodity single-machine optimization of
the grid-construction runtime. Adaptive samplers change which
grid points are evaluated; this thesis keeps the canonical grid
fixed. Heterogeneous task-scheduling work supplies the
CPU/GPU-throughput-ratio variable, but from a different workload
with different task granularity and kernels. Production ML
inference systems are closer in domain, but optimize adjacent
axes: HETERO partitions one compiled model's operation graph
across devices, automatic batching groups concurrent inference
requests along the input/request batch dimension, and PyTorch's
\texttt{functional\_call}+\texttt{vmap} pattern vectorizes a
stack of parameter states on one device. None of these directly
targets the loss-grid workload shape: one model family, many
perturbations of one checkpoint, shared data, and finite CPU,
GPU, RAM, and VRAM on a commodity machine.

